\documentclass{article}
\usepackage[utf8]{inputenc}
\usepackage[T1]{fontenc}
\usepackage{geometry}
\usepackage{xcolor}
\usepackage{tcolorbox}
\usepackage{enumitem}
\usepackage{hyperref}
\usepackage{listings}
\usepackage{tabularx}
\usepackage{multicol}
\usepackage{amssymb}

\geometry{a4paper, left=20mm, right=20mm, top=20mm, bottom=20mm}

% Custom boxes
\newtcolorbox{checklistbox}[1][]{
    colback=green!5,
    colframe=green!40!black,
    fonttitle=\bfseries,
    title=#1
}

\newtcolorbox{attackbox}[1][]{
    colback=red!5,
    colframe=red!40!black,
    fonttitle=\bfseries,
    title=#1
}

\newtcolorbox{methodologybox}[1][]{
    colback=blue!5,
    colframe=blue!40!black,
    fonttitle=\bfseries,
    title=#1
}

\title{JWT Security: \\ Checklist and Testing Methodology}
\author{Eyad Islam El-Taher}
\date{\today}

\begin{document}

\maketitle

\section{Quick Reference: JWT Attack Surface}

\begin{tabularx}{\textwidth}{|X|X|}
\hline
\textbf{Component} & \textbf{Potential Vulnerabilities} \\
\hline
Header & Algorithm tampering, \texttt{none} attack, \texttt{jwk} injection, \texttt{jku} SSRF, \texttt{kid} path traversal \\
\hline
Payload & Privilege escalation, information disclosure, missing expiration \\
\hline
Signature & Weak secret brute-force, algorithm confusion, missing verification \\
\hline
\end{tabularx}

\section{JWT Testing Methodology}

\subsection{Phase 1: Reconnaissance}

\begin{methodologybox}[Information Gathering]
\begin{enumerate}
    \item Identify where JWTs are used (cookies, Authorization headers, URL parameters)
    \item Capture multiple JWTs from different user roles (anonymous, low-privilege, admin)
    \item Check for JWT exposure in:
    \begin{itemize}
        \item \texttt{/jwks.json} or \texttt{/.well-known/jwks.json}
        \item Client-side JavaScript files
        \item API responses
        \item Error messages
    \end{itemize}
    \item Identify the algorithm used (RS256, HS256, ES256, etc.)
\end{enumerate}
\end{methodologybox}

\subsection{Phase 2: Structure Analysis}

\begin{methodologybox}[Decode and Examine]
\begin{enumerate}
    \item Decode the JWT using tools like \texttt{jwt.io} or Burp's JWT Editor
    \item Analyze header parameters:
    \begin{itemize}
        \item \texttt{alg} - algorithm
        \item \texttt{typ} - type
        \item \texttt{kid} - key ID
        \item \texttt{jwk} - embedded key
        \item \texttt{jku} - JWK set URL
        \item \texttt{x5c} - X.509 certificate
        \item \texttt{cty} - content type
    \end{itemize}
    \item Analyze payload claims:
    \begin{itemize}
        \item \texttt{sub} - subject (user)
        \item \texttt{role/group/admin} - authorization claims
        \item \texttt{exp} - expiration
        \item \texttt{iat} - issued at
        \item \texttt{aud} - audience
    \end{itemize}
\end{enumerate}
\end{methodologybox}

\subsection{Phase 3: Signature Verification Testing}

\begin{methodologybox}[Bypass Attempts]
\begin{enumerate}
    \item \textbf{Remove signature:} Delete signature part, keep trailing dot
    \item \textbf{Modify payload:} Change claims without fixing signature
    \item \textbf{Test algorithm confusion:} Change \texttt{alg} from RS256 to HS256
    \item \textbf{Test \texttt{none} algorithm:} Try various cases (none, None, NONE, none with spaces)
    \item \textbf{Null byte injection:} Add null bytes to signature or headers
\end{enumerate}
\end{methodologybox}

\subsection{Phase 4: Key-Related Attacks}

\begin{methodologybox}[Key Manipulation]
\begin{enumerate}
    \item \textbf{Brute-force weak secrets} (for HS256):
    \begin{verbatim}
    hashcat -a 0 -m 16500 <jwt> <wordlist>
    \end{verbatim}
    
    \item \textbf{jwk injection:} Embed your own public key
    \begin{verbatim}
    jwt_tool <jwt> -X i -I -pc sub -pv administrator
    \end{verbatim}
    
    \item \textbf{jku injection:} Host your JWKS and reference it
    \begin{verbatim}
    jwt_tool <jwt> -X s -ju http://attacker.com/jwks.json
    \end{verbatim}
    
    \item \textbf{kid path traversal:} Point to predictable files
    \begin{verbatim}
    "kid": "../../../../../../../dev/null"
    \end{verbatim}
\end{enumerate}
\end{methodologybox}

\subsection{Phase 5: Algorithm Confusion Attack Workflow}

\begin{methodologybox}[RS256 → HS256 Conversion]
\textbf{If public key is exposed:}
\begin{enumerate}
    \item Get public key from \texttt{/jwks.json}
    \item Convert to PEM format
    \item Base64-encode the PEM
    \item Create symmetric key with encoded PEM as \texttt{k}
    \item Change \texttt{alg} to HS256
    \item Sign with the symmetric key
\end{enumerate}

\textbf{If public key is not exposed:}
\begin{enumerate}
    \item Collect two different JWTs from the server
    \item Use sig2n tool to derive the public key:
    \begin{verbatim}
    docker run --rm -it portswigger/sig2n <token1> <token2>
    \end{verbatim}
    \item Identify correct key by testing forged JWTs
    \item Follow the same steps as above
\end{enumerate}
\end{methodologybox}

\subsection{Phase 6: Privilege Escalation Testing}

\begin{methodologybox}[Claim Manipulation]
Try modifying these common claims:
\begin{itemize}
    \item \texttt{sub: user123 → admin}
    \item \texttt{username: john → administrator}
    \item \texttt{role: user → admin, administrator, superuser}
    \item \texttt{admin: false → true}
    \item \texttt{isAdmin: 0 → 1}
    \item \texttt{group: users → administrators}
    \item \texttt{email: user@example.com → admin@example.com}
\end{itemize}
\end{methodologybox}

\section{Complete Security Checklist}
\begin{checklistbox}[JWT Security Assessment Checklist]
\subsection*{Header Validation}
\begin{itemize}
    \item [$\square$] Test \texttt{alg: none} with various cases
    \item [$\square$] Test algorithm switching (RS256 → HS256)
    \item [$\square$] Check if \texttt{jwk} parameter is accepted
    \item [$\square$] Check if \texttt{jku} parameter is accepted and validate URL fetching
    \item [$\square$] Test \texttt{kid} for path traversal (../)
    \item [$\square$] Test \texttt{kid} for SQL injection
    \item [$\square$] Test \texttt{x5c} certificate injection
    \item [$\square$] Test \texttt{cty} content type switching (XML, Java serialized)
\end{itemize}

\subsection*{Signature Verification}
\begin{itemize}
    \item [$\square$] Test if signature is verified at all
    \item [$\square$] Test if modified tokens are accepted
    \item [$\square$] Test if tokens with removed signature are accepted
    \item [$\square$] Test if tokens with modified signature are accepted
    \item [$\square$] Check for timing attacks in verification
\end{itemize}

\subsection*{Secret Key Strength}
\begin{itemize}
    \item [$\square$] Brute-force HS256 secrets using common wordlists
    \item [$\square$] Check for default or weak keys
    \item [$\square$] Test if key can be derived from multiple tokens
\end{itemize}
\end{checklistbox}

\begin{checklistbox}
\subsection*{Payload Claims}
\begin{itemize}
    \item [$\square$] Check for missing expiration (\texttt{exp})
    \item [$\square$] Test privilege escalation by modifying claims
    \item [$\square$] Check for sensitive data in payload
    \item [$\square$] Test audience (\texttt{aud}) validation
    \item [$\square$] Test issuer (\texttt{iss}) validation
\end{itemize}

\subsection*{Transmission Security}
\begin{itemize}
    \item [$\square$] Check if tokens are sent over HTTP (not HTTPS)
    \item [$\square$] Check if tokens appear in URLs (log exposure)
    \item [$\square$] Check for token leakage in referrer headers
\end{itemize}

\subsection*{Token Management}
\begin{itemize}
    \item [$\square$] Test if tokens can be revoked (logout)
    \item [$\square$] Test token reuse after logout
    \item [$\square$] Test token expiration enforcement
    \item [$\square$] Test refresh token mechanisms
\end{itemize}
\end{checklistbox}

\section{Tools Reference}

\begin{tabularx}{\textwidth}{|l|X|}
\hline
\textbf{Tool} & \textbf{Purpose} \\
\hline
Burp Suite + JWT Editor & Intercept, modify, and sign JWTs \\
\hline
jwt\_tool & Comprehensive JWT testing \\
\hline
Hashcat (mode 16500) & Brute-force HS256 secrets \\
\hline
sig2n (Docker) & Derive public keys from tokens \\
\hline
jwt.io & Quick decode and debug \\
\hline
\end{tabularx}


\end{document}